\section{Grundlagen und verwandte Arbeiten}\label{section:Grundlagen und verwandte Arbeiten}

\subsection{Verwandte Arbeiten}

Das untersuchte Feld der constraint-basierten Objektplatzierung in HAR greift in viele Bereiche der Informatik und Human-Computer Interaction (HCI) ein und kann Erkenntnisse aus verwandten Bereichen übernehmen.

Zwar wird die untersuchte Kombination des multidisziplinären Themas in bisheriger Forschung kaum betrachtet, ähnliche Probleme werden jedoch auch in Bereichen wie Game Design, Building Information Modeling (BIM) und Computer-aided Design (CAD) aufgearbeitet.

Jeder der Einzelaspekte und Anwendungsbereichen beinhaltet hierbei eine große Bandbreite an existierender Forschung. Aus diesem Grund fokussieren sich die nachfolgenden Abschnitte primär auf Werke, welche stark relevant für die Implementierung des Zielsystems sind.

\subsubsection{Constraint-basierte Platzierung}

\textcite{stuerzlingerValueConstraints3D2011a} demonstrieren in ihrer Arbeit den Mehrwert von Constraints für 3D Interfaces, um die Nutzerfreundlichkeit zu verbessern. Als Grundsatz stellen die Autoren hierfür auf, „humans are not necessarily as proficient in full 6DOF manipulation tasks as many believe“ \parencite[S. 9]{stuerzlingerValueConstraints3D2011a}. Die Möglichkeit einer uneingeschränkten Veränderung in 3D wird hierbei in Aufgaben der realen Welt kaum genutzt. So findet etwa das Laufen und Schauen in einem Flugzeug primär in vier Dimensionen statt, und weitere Dimensionen wie die Neigung des Kopfes wird nicht genutzt. Die Autoren stellen weiter auf, dass Aufgaben, welche alle DOF benötigen, nur mit speziellem, langjährigem Training möglich sind - darunter die Aufgabe von Astronauten.

Die Arbeit stellt hierfür zehn Richtlinien zur Einschränkung von 3D-Interfaces auf. Hierunter zählt unter anderem, dass nur sichtbare und keine verdeckten Objekte manipuliert werden können, oder dass die Rotation von Objekten auf logische Bewegungen reduziert wird. Als Beispiel wird hierzu ein Stuhl genannt, welcher eine feste Ausrichtung nach „oben“ erhält, da die Rotation um diese Achse sonst nicht sinnvoll ist.

\medskip

\textcite{vassionInvestigatingConstrainedObjects} untersucht die Implementierung von Constraints innerhalb von Objekten in AR. Hierfür untersucht die Arbeit, wie reale Objekte, wie Hände, mit virtuellen Objekten interagieren können. Die Constraints werden hierbei direkt in die virtuellen Objekte anhand vorgegebener Constraint-typen integriert.

Das Ergebnis der Arbeit liefert eine AR-Applikation, ähnlich wie die IKEA App, mit welcher virtuelle Möbel in die reale Umgebung platziert werden können. Der Unterschied zu bisherigen Lösungen ist es jedoch, dass mit den virtuellen Möbeln direkt interagiert werden kann, um beispielsweise eine virtuelle Schranktür mit den eigenen Händen zu öffnen.

Das System wurde hierfür mittels der Unity Game Engine, der Vuforia AR Integration und eigenen Skripts für die Implementierung der Constraints umgesetzt.

\medskip

\subsubsection{Syntax zur Constraint-Definition}

Die Definition und Auswertung von Regeln zur Platzierung von Objekten findet in diverser Kreativsoftware Anwendung. Die Speicherung und Auswertung der Regeln findet in unterschiedlichen eigenen Sprachen und Syntax statt. Die verschiedenen Ansätze werden in diesem Kapitel jeweils kurz beleuchtet, eine genauere Analyse und Klassifizierung möglicher Syntax findet stattdessen in Kapitel \ref{section:Konzeptualisierung eines Constraint-Systems} statt.

Thematisch nah zum Zielsystem sind hierbei besonders Forschungen im Bereich der Architektur-Planung, des Building Information Modeling (BIM) und der Städte-Kataster.

Im Bereich des BIM werden detaillierte Modelle von Neubauten erstellt, um Vorgänge von der Planung bis zur fertigen Bereitstellung zu modellieren \parencite{tuevsuedagWasIstBIM2025}. Städte-Kataster dagegen sind das Verzeichnis über Grundstücks-Besitzer, welche von den jeweiligen Behörden zur Verwaltung geführt werden \textcite{buettnerWasIstKataster2021}. In beiden Bereichen gibt es die Anwendung, geltende Gesetze in strikt umgesetzte Regeln umzuwandeln - etwa zu Mindestabständen zwischen Häusern.

\medskip

\textcite{kalogianniINTERLISLanguageModelling2017} beschreibt das \textit{INTERLIS}-System welches in der Schweiz zur Führung des Landes-Kataster genutzt wird. Es enthält neben der Geometrie auch Informationen zu Attributen, sozialen sowie rechtlichen Semantiken bestimmter Grundstücke in 3D.

Die Arbeit stellt zunächst die Object Constraint Language (OCL) vor, welche als Erweiterung von UML fungiert. Diese ermöglicht es, zwischen vorher entwickelten UML-Modellen Regeln aufzustellen. Als zentraler Nachteil wird jedoch genannt, dass die in OCL definierten Regeln primär in der Phase des Systems Engineering Anwendung finden und nicht direkt in Programmiercode umgewandelt oder ausgewertet werden können.

\begin{lstlisting}[frame=htrbl, caption={INTERLIS Constraint Language Beispiel}, label={lst:interlis}, language=php]
CLASS Bench EXTENDS Furniture =
  location: MANDATORY SURFACE WITH (STRAIGHTS);
  surfaceType: MANDATORY (grass, asphalt, concrete, sand);

MANDATORY CONSTRAINT
  surfaceType == grass OR surfaceType == sand;
END Bench;
\end{lstlisting}


Constraints in \textit{INTERLIS} werden stattdessen über die eigene \textit{INTERLIS constraint language} definiert. Listing \ref{lst:interlis} bietet ein Beispiel für eine vollständige Constraint-Definition, abgeleitet für den Anwendungsfall dieser Arbeit. Der Code definitiert zunächst die Klasse „Bench“, welches von der Grundklasse „Furniture“ ableitet. Für die Klasse werden die Attribute der Position und Oberfläche des Untergrunds definiert. Über das Stichwort „Mandatory Constraint“ kann schließlich eine Regel definiert werden, dass die Bank nur auf Gras oder Sand gestellt werden darf.

\medskip

\textcite{sydoraRulebasedComplianceChecking2020} nehmen als Grundlage für die Definition stattdessen mathematische Ausdrücke. 

\begin{lstlisting}[frame=htrbl, caption={BIM Constraint language \parencite{sydoraRulebasedComplianceChecking2020}}, label={lst:bim-constraint}, language=php, literate={∈}{{$\in$}}1]
ALL Obj0 ∈ IfcFurnishing {FunctionOfObj(Obj0) = "Refrigerator"};
ANY Obj1 ∈ IfcCounterTop;

(IsNextTo(Obj1, Obj0) = TRUE)
\end{lstlisting}

Listing \ref{lst:bim-constraint} zeigt ein verkürztes Beispiel für ein Ausdruck in der definierten Sprache. Das Beispiel definiert, dass alle Möbel, welche die Funktion des Kühlschranks haben (\textit{Obj0}), eine Arbeitsplatte besitzen müssen (\textit{Obj1}), welche direkt angrenzend ist. Hierfür werden zusätzlich Stichwörter wie \textit{ALL} und \textit{ANY} verwendet, um zu definieren, dass \textit{jeder} Kühlschrank \textit{irgendeine} Arbeitsplatte in der Szene besitzen muss.

Für die Identifikation bestimmer Elemente und Eigenschaften nutzt das System, ähnlich wie das Tagging von Objekten, eigene Objekt-Attribute. Für die Erstellung neuer Regeln wird ein einfacher Regel-Editor, ähnlich einem Code-Editor, im Paper konzeptioniert. 

Trotz dieser mathematischen Korrektheit der Regeln weisen die Autoren darauf hin, dass definierte Regeln - egal über welches System - immer mehrdeutig sein können. Als Beispiel wird hierzu eine scheinbar einfache Regel zum Abstand zwischen zwei Gegenständen genannt: Dieser Abstand könnte jedoch im Code als Abstand zwischen den Mittelpunkten der Objekte, als Abstand auf einer 2D-Plane, auf Grundlage einer rechteckigen Bounding Box oder anderen Berechnungen berechnet werden und je nachdem ein anderes Ergebnis haben.

\medskip

\textcite{schwabeBIMApplicationsRuleBased2016a}

\medskip

//solibri


\subsubsection{Kartenintegration}

\subsubsection{Nutzer-Feedback}